\section{Training and Sampling Hyperparameters}
\label{app:training-sampling-hyperparameters}

We summarize the hyperparameters used for \model{}-M and \model{}-L.
Unless noted, all settings are shared between the two model sizes.

\textbf{Architecture.}
\Cref{tab:ours-architecture-hyperparameters} specifies model capacity and
backbone configuration; the two models differ only in width, head count, and
parameter count.

\begin{table}[H]
\centering
\caption{\textbf{Architecture hyperparameters.}
Model capacity and backbone configuration for \ours{}-M and \ours{}-L.}
\label{tab:ours-architecture-hyperparameters}
\begin{tabular}{@{}lrr@{}}
\toprule
Parameter & \ours{}-M & \ours{}-L \\
\midrule
\multicolumn{3}{@{}l}{\textbf{\emph{Capacity}}} \\
\quad Parameters       & 88M & 187M \\
\quad Single width     & 512 & 768 \\
\quad Pair width       & 128 & 192 \\
\quad Attention heads  & 8   & 12 \\
\addlinespace
\multicolumn{3}{@{}l}{\textbf{\emph{Backbone}}} \\
\quad Pairmixer blocks & 4 & 4 \\
\quad DiT blocks       & 16 & 16 \\
\quad MLP ratio        & 4 & 4 \\
\quad QK normalization & On & On \\
\quad Feed-forward network & SwiGLU & SwiGLU \\
\quad Dropout          & 0 & 0 \\
\bottomrule
\end{tabular}
\end{table}

\textbf{Training.}
\Cref{tab:ours-training-hyperparameters} summarizes the flow-matching,
optimization, conditioning, and checkpoint-selection settings.

\begin{table}[H]
\centering
\caption{\textbf{Training hyperparameters.}
Flow-matching, optimization, conditioning, and model-selection settings.
All settings are shared between \ours{}-M and \ours{}-L unless noted.}
\label{tab:ours-training-hyperparameters}
\begin{tabular}{@{}>{\raggedright\arraybackslash}p{0.39\columnwidth}
                    >{\raggedright\arraybackslash}p{0.55\columnwidth}@{}}
\toprule
Parameter & Value \\
\midrule
\multicolumn{2}{@{}l}{\textbf{\emph{Flow matching}}} \\
\quad Priors & Centered Gaussian fpr $\mathbf{X}$; standard Gaussian for $\boldsymbol{\ell}$ \\
\quad Interpolant & Linear ($x_t=(1-t)x_0+tx_1$) \\
\quad Time distribution &
$0.98\,\operatorname{Beta}(1.9,1.0)+0.02\,\operatorname{Uniform}(0,1)$ \\
\quad Objective & $L^1$ coordinate--lattice loss \\
\quad Loss weights & Coordinates 10; lattice 1 \\
\addlinespace
\multicolumn{2}{@{}l}{\textbf{\emph{Optimization}}} \\
\quad Optimizer & Muon for hidden matrices; AdamW otherwise \\
\quad Learning rate & $10^{-4}$ \\
\quad Weight decay & 0 (\ours{}-M); $10^{-2}$ (\ours{}-L) \\
\quad Muon settings & Momentum 0.95; Nesterov; 5 Newton--Schulz steps \\
\quad AdamW settings & $\beta_1=0.9$, $\beta_2=0.999$, $\epsilon=10^{-8}$ \\
\quad Schedule & 10{,}000-step linear warmup from zero \\
\quad Gradient clipping & Global norm, 1.0 \\
\quad EMA decay & 0.9999 \\
\addlinespace
\multicolumn{2}{@{}l}{\textbf{\emph{Conditioning and augmentation}}} \\
\quad Condition dropout & Template 0.5; stereochemistry 0.5; space group 0.9 \\
\quad Augmentation & Crystal translation; template rotation and translation \\
\addlinespace
\multicolumn{2}{@{}l}{\textbf{\emph{Runtime and selection}}} \\
\quad Precision & bfloat16 mixed precision \\
\quad Checkpoint interval & Every 50 epochs \\
\quad Checkpoint selection & Highest validation $\mathrm{Sol}_C$ \\
\bottomrule
\end{tabular}
\end{table}

\textbf{Curriculum and batching.}
We use two-stage training with fixed-shape batches to enable
\texttt{torch.compile}. Each structure is assigned to the smallest
atom-count bucket that contains it and padded to that bucket size.

\begin{table}[H]
\centering
\caption{\textbf{Two-stage training curriculum and batching.}
Atom-count buckets and batch sizes used for fixed-shape training across eight GPUs.}
\label{tab:ours-curriculum-batches}
\begin{tabular}{@{}>{\raggedright\arraybackslash}p{0.39\columnwidth}
                    >{\raggedright\arraybackslash}p{0.55\columnwidth}@{}}
\toprule
Parameter & Value \\
\midrule
\multicolumn{2}{@{}l}{\textbf{\emph{Stage 1}}} \\
\quad Maximum atoms & 300 \\
\quad Bucket caps & 32, 64, 96, 128, 160, 192, 224, 256, 288, 300 \\
\quad Per-GPU batch & 16 for every bucket \\
\quad Global batch & 128 \\
\addlinespace
\multicolumn{2}{@{}l}{\textbf{\emph{Stage 2}}} \\
\quad Maximum atoms & 512 \\
\quad Additional bucket caps & 320, 352, 384, 416, 448, 480, 512 \\
\quad Per-GPU batch & 16 ($\leq300$), 14 (320), 11 (352), 9 (384),
8 (416), 7 (448), 6 (480), 5 (512) \\
\quad Global batch & 128, 112, 88, 72, 64, 56, 48, and 40, respectively \\
\bottomrule
\end{tabular}
\end{table}

\textbf{Sampling.}
Both model sizes use the same sampling and conditioning configuration for
benchmark generation.

\begin{table}[H]
\centering
\caption{\textbf{Sampling hyperparameters.}
Sampling and conditioning configuration used for benchmark generation.}
\label{tab:ours-sampling-hyperparameters}
\begin{tabular}{@{}>{\raggedright\arraybackslash}p{0.39\columnwidth}
                    >{\raggedright\arraybackslash}p{0.55\columnwidth}@{}}
\toprule
Parameter & Value \\
\midrule
\multicolumn{2}{@{}l}{\textbf{\emph{Numerical solver}}} \\
\quad Solver & EDM--Heun with stochastic churn \\
\quad Noise schedule & Karras \\
\quad Sampling steps & 200 \\
\quad Sigma range & $\sigma_{\min}=0.002$, $\sigma_{\max}=80$ \\
\quad Schedule exponent & $\rho=7$ \\
\quad Churn settings & $S_{\mathrm{churn}}=60$, $S_{\min}=0$,
$S_{\max}=999$, $S_{\mathrm{noise}}=1.003$ \\
\addlinespace
\multicolumn{2}{@{}l}{\textbf{\emph{Generation protocol}}} \\
\quad Weights & EMA \\
\quad Conditioning & Template and available stereochemistry on; space group off \\
\quad Template augmentation & Random rotation and translation \\
\quad Candidates per target & 1{,}000 \\
\quad Random seed & 42 \\
\bottomrule
\end{tabular}
\end{table}
